\chapter{Coccidioidomycosis}
\index{Coccidioidomycosis}

\section*{Definition and Overview}
Coccidioidomycosis, commonly known as Valley fever, is a systemic fungal infection caused by the dimorphic soil-dwelling fungi \textit{Coccidioides immitis} or \textit{Coccidioides posadasii}. The disease is endemic to arid and semi-arid regions of the Western Hemisphere, particularly the southwestern United States (Arizona, California), parts of Mexico, and Central and South America. Infection occurs via the inhalation of airborne arthroconidia (fungal spores) dislodged from the soil. The spectrum of disease ranges from asymptomatic or mild self-limiting respiratory tract infections to severe, progressive pulmonary disease or life-threatening disseminated infection affecting the skin, bones, joints, and meninges.

\section*{Epidemiology}
Coccidioidomycosis is highly prevalent in endemic regions. In the United States, tens of thousands of cases are reported annually, with Arizona and California accounting for the vast majority. The disease is not endemic to some countries, but is clinically significant in travellers returning from endemic areas. Anyone visiting these regions can acquire the infection, but certain groups are at higher risk for severe or disseminated disease, including older adults, pregnant women (particularly in the third trimester), and immunocompromised individuals (such as those with HIV/AIDS, organ transplant recipients, or patients on anti-TNF therapy). People of African or Filipino descent also have a disproportionately higher risk of dissemination for reasons that are not fully understood.

\section*{Aetiology and Pathophysiology}
Coccidioidomycosis is caused by the dimorphic fungi of the genus \textit{Coccidioides}:

\begin{itemize}
    \item \textbf{Dimorphic lifecycle:} in the soil, the fungus grows as a mold, forming barrel-shaped arthroconidia. Once inhaled, the arthroconidia transform in the warm, moist environment of the human lung into thick-walled spherules.
    \item \textbf{Spherule rupture:} as spherules mature, they fill with thousands of small endospores. The spherules eventually rupture, releasing endospores into the surrounding lung tissue, each of which can form new spherules, propagating the infection.
    \item \textbf{Immune response:} the host immune system mounts a granulomatous and suppurative inflammatory response to contain the infection. If cell-mediated immunity is impaired, the fungus can escape the lungs and disseminate hematogenously to other organs.
\end{itemize}

\section*{Clinical Features}
Approximately 60\% of infected individuals are asymptomatic. In symptomatic patients, features present 1 to 3 weeks after inhalation:

\begin{itemize}
    \item Primary pulmonary infection: presenting as a flu-like illness with fever, cough, pleuritic chest pain, sore throat, headache, fatigue, and muscle aches
    \item Dermatological signs: erythema nodosum (painful red nodules on the shins) or erythema multiforme, which are immunologic hypersensitivity reactions indicating a robust host immune response and generally favorable prognosis
    \item Disseminated disease: presenting with chronic skin lesions (papules, pustules, or ulcers), osteomyelitis, septic arthritis, and chronic meningitis (manifesting as persistent headache, stiff neck, and confusion)
\end{itemize}

\section*{Differential Diagnosis}
Coccidioidomycosis can mimic other respiratory and infectious conditions:

\begin{itemize}
    \item \textbf{Bacterial pneumonia:} presenting with similar fever, cough, and infiltrate on chest X-ray, but responding to antibacterial therapy.
    \item \textbf{Tuberculosis (TB):} presenting with chronic cough, night sweats, weight loss, and upper lobe cavitary lesions, requiring acid-fast bacilli smear and culture to differentiate.
    \item \textbf{Other fungal infections:} such as histoplasmosis or blastomycosis, which have different endemic geographic distributions and microscopic appearances.
    \item \textbf{Lung cancer:} cavitary or nodular pulmonary lesions on imaging can be mistaken for malignancy.
    \item \textbf{Meningitis of other causes:} bacterial, viral, or tuberculous meningitis must be excluded through cerebrospinal fluid analysis.
\end{itemize}

\section*{When to Seek Medical Care}
Anyone who has travelled to or resided in an endemic area and develops persistent flu-like symptoms, chronic cough, joint swelling, or skin lesions should consult a doctor.

\textbf{Contact your local emergency services for:}
\begin{itemize}
    \item Severe shortness of breath or difficulty breathing
    \item Sudden, severe chest pain
    \item Signs of meningitis: severe, intractable headache, stiff neck, high fever, and sensitivity to light
    \item Confusion, altered consciousness, or seizures
\end{itemize}

\section*{Diagnosis and Investigations}
\begin{itemize}
    \item \textbf{Serology:} detection of IgM and IgG antibodies using enzyme immunoassay (EIA), immunodiffusion, or complement fixation. IgG titres are useful for monitoring disease severity and response to treatment.
    \item \textbf{Microscopic examination:} identification of characteristic spherules in sputum, bronchoalveolar lavage fluid, or tissue biopsy specimens using special stains (e.g., KOH preparation or GMS stain).
    \item \textbf{Fungal culture:} isolation of \textit{Coccidioides} from clinical specimens. Note: because the mold phase is highly infectious, laboratory personnel must handle specimens with extreme caution under Biosafety Level 3 conditions.
    \item \textbf{Imaging:} chest X-ray or CT scans showing lobar infiltrates, nodules, thin-walled cavities, or hilar lymphadenopathy.
    \item \textbf{Lumbar puncture:} essential if coccidioidal meningitis is suspected, showing lymphocytic pleocytosis, elevated protein, and low glucose in cerebrospinal fluid.
\end{itemize}

\section*{Management}
Mild, uncomplicated cases in low-risk individuals may require only close observation. Treatment is indicated for severe or progressive disease, and all patients with risk factors for dissemination:

\begin{itemize}
    \item \textbf{Antifungal therapy:} oral azoles, primarily fluconazole or itraconazole, are first-line agents for pulmonary and non-meningeal disseminated disease.
    \item \textbf{Amphotericin B:} reserved for severe, rapidly progressive infections or in pregnant patients where azoles are teratogenic.
    \item \textbf{Meningitis management:} requires lifelong treatment with high-dose oral fluconazole to prevent relapse.
    \item \textbf{Surgical intervention:} may be needed for joint debridement, drainage of abscesses, or resection of persistent, bleeding pulmonary cavities.
\end{itemize}

\section*{Complications}
\begin{itemize}
    \item Chronic progressive pulmonary disease with persistent cavitary lesions, which may bleed (haemoptysis) or rupture into the pleural space
    \item Disseminated disease leading to severe bone and joint destruction
    \item Coccidioidal meningitis, which carries high mortality and morbidity if untreated and can cause hydrocephalus
    \item Adverse effects of prolonged antifungal therapy (e.g., hepatotoxicity)
\end{itemize}

\section*{Prognosis and Outcomes}
The prognosis for primary pulmonary coccidioidomycosis in healthy individuals is excellent, with symptoms resolving fully over weeks to months without antifungal treatment. For patients requiring treatment, oral azole therapy is highly effective, though recovery can be slow. The prognosis for disseminated disease is more guarded, and coccidioidal meningitis requires lifelong medication, as discontinuation leads to a very high rate of neurological relapse.

\section*{Prevention and Self-Care}
\begin{itemize}
    \item Avoid exposure to dust and soil in endemic areas, particularly during dust storms, construction, or excavation
    \item Stay indoors during windy conditions and keep windows closed in endemic regions
    \item Consider wearing a high-efficiency particulate air (HEPA) mask (such as an N95 respirator) if outdoor activity in dusty environments is unavoidable in endemic zones
    \item Wet down soil prior to digging to prevent spores from becoming airborne
    \item Learn to recognize symptoms early, especially if you are in a high-risk group (e.g., pregnant or immunocompromised)
\end{itemize}

\section*{Sources and Further Reading}
The following reputable organisations provide further information for healthcare students and patients seeking reliable, current advice about coccidioidomycosis.

\begin{itemize}
    \item Centers for Disease Control and Prevention (CDC). \textit{Valley fever (coccidioidomycosis) information and resources.} https://www.cdc.gov/fungal/diseases/coccidioidomycosis/
    \item Arizona Department of Health Services. \textit{Endemic updates, clinical guidelines, and patient fact sheets.} https://www.azdhs.gov/
    \item Mayo Clinic. \textit{Valley fever: symptoms, causes, diagnosis, and treatment.} https://www.mayoclinic.org/
    \item MSD Manual. \textit{Coccidioidomycosis: professional and patient reference.} https://www.msdmanuals.com/
    \item World Health Organization. \textit{Information on endemic fungal infections and public health.} https://www.who.int/
\end{itemize}
